% Original analytic teaching diagram, 2026-11-12.
% Two dice, 36 equally likely ordered outcomes (i,j). (a) each cell shows S=i+j; the atom {S=7} is filled,
% the red frame is the event {first die = 1}, which cuts across atoms of sigma(S).
% (b) each cell shows M=max(i,j); the atom {M=4} is filled.
\definecolor{VB}{HTML}{7188A3}\definecolor{VBF}{HTML}{E5EAF0}
\definecolor{VG}{HTML}{829C8B}\definecolor{VGF}{HTML}{E8EEE9}
\definecolor{VY}{HTML}{B4A16F}\definecolor{VYF}{HTML}{F1ECDD}
\definecolor{VR}{HTML}{A7787D}\definecolor{VRF}{HTML}{F0E5E6}
\definecolor{VP}{HTML}{9588A5}\definecolor{VPF}{HTML}{ECE8F0}
\definecolor{VI}{HTML}{3B4252}\definecolor{VM}{HTML}{6B7280}\definecolor{VL}{HTML}{4C4D4F}
\begin{tikzpicture}[x=1mm,y=1mm,>=Stealth,font=\figurefont\fontsize{8}{10}\selectfont,
lab/.style={text=VM,font=\figurefont\fontsize{7.5}{9}\selectfont},
num/.style={text=VI,font=\figurefont\fontsize{7}{8}\selectfont}]
\def\c{5.5}
% ---------- panel (a): S = i + j ----------
\begin{scope}[shift={(0,0)}]
\foreach \i in {1,...,6} { \foreach \j in {1,...,6} {
  \pgfmathtruncatemacro{\s}{\i+\j}
  \ifnum\s=7 \fill[VBF] ({(\i-1)*\c},{(\j-1)*\c}) rectangle ({\i*\c},{\j*\c}); \fi
  \draw[VL,line width=.35pt] ({(\i-1)*\c},{(\j-1)*\c}) rectangle ({\i*\c},{\j*\c});
  \node[num] at ({(\i-.5)*\c},{(\j-.5)*\c}) {\s};
}}
\draw[VR,line width=1pt] (0,0) rectangle (\c,{6*\c});
\foreach \i in {1,...,6} { \node[lab] at ({(\i-.5)*\c},-2.6) {\i}; \node[lab] at (-2.6,{(\i-.5)*\c}) {\i}; }
\node[text=VI] at ({3*\c},-7) {第一枚};
\node[text=VI,rotate=90] at (-7,{3*\c}) {第二枚};
\node[text=VI] at ({3*\c},-12) {(a)};
\end{scope}
% ---------- panel (b): M = max(i,j) ----------
\begin{scope}[shift={(52,0)}]
\foreach \i in {1,...,6} { \foreach \j in {1,...,6} {
  \pgfmathtruncatemacro{\m}{max(\i,\j)}
  \ifnum\m=4 \fill[VBF] ({(\i-1)*\c},{(\j-1)*\c}) rectangle ({\i*\c},{\j*\c}); \fi
  \draw[VL,line width=.35pt] ({(\i-1)*\c},{(\j-1)*\c}) rectangle ({\i*\c},{\j*\c});
  \node[num] at ({(\i-.5)*\c},{(\j-.5)*\c}) {\m};
}}
\foreach \i in {1,...,6} { \node[lab] at ({(\i-.5)*\c},-2.6) {\i}; \node[lab] at (-2.6,{(\i-.5)*\c}) {\i}; }
\node[text=VI] at ({3*\c},-7) {第一枚};
\node[text=VI,rotate=90] at (-7,{3*\c}) {第二枚};
\node[text=VI] at ({3*\c},-12) {(b)};
\end{scope}
\end{tikzpicture}
